% ============================================================
%  AI 产品经理手册 · 第零章第一节
%  生成式 AI 可能是人类信息革命的终局
% ============================================================
\documentclass[12pt, a4paper, openany]{book}

% ── 字体与语言 ──────────────────────────────────────────────
\usepackage{fontspec}
\usepackage{xeCJK}

\setCJKmainfont[
  BoldFont   = Heiti SC,
  ItalicFont = Kaiti SC
]{Songti SC}
\setCJKsansfont{Heiti SC}
\setmainfont{Georgia}
\setsansfont{Helvetica Neue}

% ── 页面布局 ────────────────────────────────────────────────
\usepackage[
  a4paper,
  top=2.8cm, bottom=2.8cm,
  inner=3.2cm, outer=2.4cm,
  headheight=15pt
]{geometry}

% ── 行距与段落 ──────────────────────────────────────────────
\usepackage{setspace}
\setstretch{1.6}
\usepackage{indentfirst}
\setlength{\parindent}{2em}
\usepackage{parskip}
\setlength{\parskip}{0.4em}

% ── 颜色 ────────────────────────────────────────────────────
\usepackage{xcolor}
\definecolor{AccentBlue}{RGB}{28, 69, 135}
\definecolor{AccentGold}{RGB}{180, 140, 40}
\definecolor{LightBg}{RGB}{243, 246, 251}
\definecolor{WarmGray}{RGB}{100, 95, 88}
\definecolor{TimelineGray}{RGB}{180, 180, 180}
\definecolor{HighlightYellow}{RGB}{255, 248, 220}
\definecolor{DarkText}{RGB}{30, 30, 30}

% ── 页眉页脚 ────────────────────────────────────────────────
\usepackage{fancyhdr}
\pagestyle{fancy}
\fancyhf{}
\fancyhead[LE]{\small\color{WarmGray}\itshape AI 产品经理手册}
\fancyhead[RO]{\small\color{WarmGray}\itshape 第零章\ 新起点}
\fancyfoot[C]{\small\color{WarmGray}\thepage}
\renewcommand{\headrulewidth}{0.3pt}
\renewcommand{\headrule}{\color{TimelineGray}\hrule height 0.3pt}

% ── 章节标题样式 ────────────────────────────────────────────
\usepackage{titlesec}
\titleformat{\chapter}[display]
  {\normalfont\fontsize{28}{32}\selectfont\bfseries\color{AccentBlue}}
  {\normalfont\large\color{WarmGray}第\ \thechapter\ 章}
  {10pt}
  {}
  [\vspace{6pt}\color{AccentGold}\rule{\textwidth}{1.5pt}]

\titleformat{\section}
  {\normalfont\fontsize{16}{20}\selectfont\bfseries\color{AccentBlue}}
  {\thesection}
  {0.8em}
  {}
  [\vspace{2pt}\color{AccentBlue!30}\rule{0.5\textwidth}{0.5pt}]

\titleformat{\subsection}
  {\normalfont\fontsize{13}{16}\selectfont\bfseries\color{DarkText}}
  {\thesubsection}
  {0.6em}
  {}

\titlespacing*{\chapter}{0pt}{-10pt}{30pt}
\titlespacing*{\section}{0pt}{20pt}{8pt}
\titlespacing*{\subsection}{0pt}{14pt}{4pt}

% ── 图形 & TikZ ─────────────────────────────────────────────
\usepackage{graphicx}
\usepackage{tikz}
\usepackage{pgfplots}
\pgfplotsset{compat=1.18}
\usetikzlibrary{
  arrows.meta,
  shapes.geometric,
  shapes.symbols,
  shapes.misc,
  positioning,
  fit,
  backgrounds,
  decorations.pathreplacing,
  decorations.markings,
  calc,
  shadows.blur,
  patterns
}

% ── 文本框 ──────────────────────────────────────────────────
\usepackage[most]{tcolorbox}
\tcbuselibrary{skins, breakable}

\newtcolorbox{coreview}{
  enhanced,
  breakable,
  colback=LightBg,
  colframe=AccentBlue,
  boxrule=0.8pt,
  arc=5pt,
  left=14pt, right=14pt,
  top=10pt, bottom=10pt,
  fonttitle=\bfseries\color{AccentBlue},
  title={\footnotesize 核心观点},
  attach boxed title to top left={yshift=-2mm, xshift=6mm},
  boxed title style={
    colback=AccentBlue, colframe=AccentBlue,
    arc=3pt, boxrule=0pt,
    fontupper=\small\bfseries\color{white}
  }
}

\newtcolorbox{blockquote}{
  enhanced,
  colback=HighlightYellow,
  colframe=AccentGold,
  boxrule=0pt,
  leftrule=4pt,
  arc=2pt,
  left=14pt, right=14pt,
  top=8pt, bottom=8pt,
}

\newtcolorbox{thoughtbox}{
  enhanced,
  breakable,
  colback=white,
  colframe=AccentGold,
  boxrule=1.2pt,
  arc=6pt,
  left=16pt, right=16pt,
  top=12pt, bottom=12pt,
  fonttitle=\bfseries\large\color{AccentGold},
  title={思考题},
  attach boxed title to top left={yshift=-3mm, xshift=8mm},
  boxed title style={
    colback=white, colframe=AccentGold,
    arc=3pt, boxrule=1.2pt,
    fontupper=\bfseries\color{AccentGold}
  }
}

% ── 图表标题 ────────────────────────────────────────────────
\usepackage{caption}
\captionsetup{
  font={small, color=WarmGray},
  labelfont={bf, color=AccentBlue},
  labelsep=quad,
  skip=6pt
}

% ── 表格 ────────────────────────────────────────────────────
\usepackage{booktabs}
\usepackage{array}
\usepackage{tabularx}

% ── 超链接 ──────────────────────────────────────────────────
\usepackage{hyperref}
\hypersetup{
  colorlinks=true,
  linkcolor=AccentBlue,
  urlcolor=AccentBlue,
  pdftitle={AI 产品经理手册},
  pdfauthor={AI Handbook}
}

% ── 自定义命令 ──────────────────────────────────────────────
\newcommand{\key}[1]{\textbf{\color{AccentBlue}#1}}
\newcommand{\emphasis}[1]{\textit{\color{AccentGold}#1}}
\newcommand{\data}[1]{\textbf{#1}}

% ============================================================
\begin{document}
\frontmatter

% ── 标题页 ──────────────────────────────────────────────────
\thispagestyle{empty}
\vspace*{4cm}
\begin{center}
  {\color{WarmGray}\footnotesize AI\ \ 产品经理手册\ \ ·\ \ 第零章}\\[6pt]
  \color{AccentBlue}\rule{8cm}{1pt}\\[18pt]
  {\fontsize{28}{34}\selectfont\bfseries 生成式\ AI\ 可能是\\[6pt]人类信息革命的终局}\\[20pt]
  \color{AccentBlue}\rule{8cm}{0.5pt}\\[20pt]
  {\large\color{WarmGray} 理解这次变革的量级，是做好\ AI\ 产品的前提}
\end{center}
\vfill
\begin{center}
  \color{TimelineGray}\small\itshape 2025
\end{center}
\newpage

\mainmatter
\setcounter{chapter}{0}

% ============================================================
\chapter{新起点：理解这次技术变革}
% ============================================================

\vspace{-4pt}
\begin{blockquote}
\itshape\color{WarmGray}
``人类曾经以采集食物为生，而如今他们重新要以采集信息为生。''\\
\hfill —— 马歇尔·麦克卢汉（Marshall McLuhan）
\end{blockquote}
\vspace{10pt}

如果麦克卢汉还在世，或许他会修改自己的这句话：
\emphasis{未来，人类将以创造为生。}

这不是口号，而是对一次结构性变革的判断。要理解这个判断，
我们需要先回到一个更古老的问题——信息，到底是怎么演化到今天这一步的。

% ============================================================
\section{信息的爆炸与分发的革命}
% ============================================================

\subsection{从口口相传到算法推荐}

人类历史上，信息技术的每一次跃迁，都不只是工具的升级，
而是权力结构和社会组织方式的重构。

从口口相传，到甲骨文、竹简，到印刷纸质书籍，再到电报、电台、电视、
电脑、手机——每一次媒介的变革，人类每天产生和流通的信息量，
都以数量级跳跃式增长。

\vspace{12pt}

% ── 图1：信息革命时间轴 ──────────────────────────────────────
\begin{figure}[h]
\centering
\begin{tikzpicture}[
  every node/.style={font=\small},
  milestone/.style={
    circle, draw=AccentBlue, fill=white,
    line width=1.2pt, minimum size=9pt, inner sep=0
  },
  label above/.style={
    text width=2.8cm, align=center, above=12pt,
    font=\scriptsize\color{DarkText}
  },
  label below/.style={
    text width=2.8cm, align=center, below=12pt,
    font=\scriptsize\color{WarmGray}
  },
  era/.style={
    rectangle, fill=AccentBlue!8, draw=AccentBlue!20,
    rounded corners=3pt,
    font=\scriptsize\bfseries\color{AccentBlue}
  }
]

% 主轴
\draw[-{Stealth[length=6pt]}, line width=1.2pt, color=AccentBlue!50]
  (0,0) -- (14.8,0);

% 刻度点
\foreach \x/\lbl/\yr/\pos in {
  0.6/口口相传/远古/above,
  2.0/甲骨文竹简/3000 BC/below,
  3.6/活字印刷/1450 AD/above,
  5.2/电报电话/1850s/below,
  6.8/广播电视/1920s/above,
  8.4/互联网/1990s/below,
  10.0/移动推荐/2010s/above,
  11.8/生成式AI/2020s/below
}{
  \node[milestone] at (\x, 0) {};
  \ifnum\pdfstrcmp{\pos}{above}=0
    \draw[TimelineGray, dashed, thin] (\x, 0) -- (\x, 0.9);
    \node[label above] at (\x, 0) {\lbl\\\color{WarmGray}\itshape\yr};
  \else
    \draw[TimelineGray, dashed, thin] (\x, 0) -- (\x, -0.9);
    \node[label below] at (\x, 0) {\lbl\\\color{WarmGray}\itshape\yr};
  \fi
}

% 最后一个节点特殊标注
\node[circle, draw=AccentGold, fill=AccentGold!20,
      line width=1.8pt, minimum size=11pt, inner sep=0] at (11.8, 0) {};

% 分段颜色渐变标注
\node[era, minimum width=3.8cm] at (2.8, -2.2) {手工记录时代};
\node[era, minimum width=3.6cm] at (6.0, 2.2)  {电气广播时代};
\node[era, minimum width=3.4cm] at (9.2, -2.2) {数字互联网时代};
\node[era, fill=AccentGold!15, draw=AccentGold!40,
      minimum width=2.6cm] at (11.8, 2.2) {AI时代};

\end{tikzpicture}
\caption{人类信息媒介演进时间轴：从口耳相传到生成式AI，每一次跃迁都重塑了
信息的生产权与传播权}
\label{fig:timeline}
\end{figure}

\vspace{6pt}

面对指数级增长的信息洪流，人类也在持续发明新的分发方式。
图书馆的分类编目系统，是对有限知识的有序归档；
搜索引擎的出现，让"平等找到所求"第一次成为可能——
贵州山区的孩子与北京的学生，在理论上可以检索到同等质量的信息。
个性化推荐算法更进一步，它不再等待你去"找"，
而是主动把信息推送到你面前。

这三代分发技术，本质上都在做同一件事：
\key{帮助人类对抗信息过载，提升获取有价值信息的效率}——提高信噪比。

\begin{coreview}
从图书馆到搜索引擎，再到今天的个性化推荐，
信息\textbf{消费}的门槛一降再降。
一个偏远地区的孩子，今天也能实时获取到最前沿的知识。
这是分发技术带来的\textbf{消费侧平权}，
它已基本走到了技术演进的成熟阶段。
\end{coreview}

% ============================================================
\section{生产：从未真正平权}
% ============================================================

当我们把信息的生产与消费对比来看，会发现一个长期被忽视的结构性不对称：
\key{消费侧已相对平权，生产侧从未真正平权。}

\subsection{谁有资格"说话"：生产权的历史}

信息的生产权，在人类历史的绝大多数时间里，都是少数人的特权。

在古代，信息的生产者是祭司和史官——他们掌握文字，
垄断着对自然与历史的解释权；消费者是当权者，普通民众甚至不在分发链路之内。
随着科举和私塾的兴起，"读书人"成为一种职业身份，
代写文书、为人读信是专业能力，文字开始承载学习与交流的价值。
进入电台电视时代，信息的生产依赖昂贵的技术设备和播出资质，
生产权高度集中在国家与媒体机构手中。

互联网时代似乎打破了这一格局——任何人都可以在社交平台发布内容。
然而现实数据却令人警醒：\data{在今天的短视频平台上，
每天真正持续创作内容的用户，不足总用户的 1.5\%}。

\vspace{14pt}

% ── 图2：生产权的历史金字塔 ─────────────────────────────────
\begin{figure}[h]
\centering
\begin{tikzpicture}[
  every node/.style={font=\small\color{white}},
  scale=0.88
]

\def\pw{10}  % 金字塔最大宽度

% 各层：从底到顶
% Layer 5: AI时代（未来，最宽/底层翻转）
\fill[AccentGold!70]
  ({-\pw/2*1.0}, -5.2) -- ({\pw/2*1.0}, -5.2)
  -- ({\pw/2*0.85}, -4.2) -- ({-\pw/2*0.85}, -4.2) -- cycle;
\node[font=\small\bfseries\color{DarkText}] at (0, -4.7)
  {AI 时代：\textbf{所有人}皆可创作};

% Layer 4: 互联网时代
\fill[AccentBlue!55]
  ({-\pw/2*0.85}, -4.2) -- ({\pw/2*0.85}, -4.2)
  -- ({\pw/2*0.6}, -3.1) -- ({-\pw/2*0.6}, -3.1) -- cycle;
\node[font=\scriptsize\color{white}] at (0, -3.65)
  {互联网时代\ \ 创作者\ \textasciitilde1.5\%};

% Layer 3: 广电时代
\fill[AccentBlue!72]
  ({-\pw/2*0.6}, -3.1) -- ({\pw/2*0.6}, -3.1)
  -- ({\pw/2*0.38}, -2.1) -- ({-\pw/2*0.38}, -2.1) -- cycle;
\node[font=\scriptsize\color{white}] at (0, -2.6)
  {广电时代\ \ 专业机构\ \textless0.1\%};

% Layer 2: 印刷文人时代
\fill[AccentBlue!85]
  ({-\pw/2*0.38}, -2.1) -- ({\pw/2*0.38}, -2.1)
  -- ({\pw/2*0.22}, -1.2) -- ({-\pw/2*0.22}, -1.2) -- cycle;
\node[font=\scriptsize\color{white}] at (0, -1.65)
  {文人书写时代\ \ 读书人\ \textasciitilde5\%};

% Layer 1: 古代祭祀时代（顶端）
\fill[AccentBlue]
  ({-\pw/2*0.22}, -1.2) -- ({\pw/2*0.22}, -1.2)
  -- (0, -0.2) -- cycle;
\node[font=\scriptsize\color{white}] at (0, -0.8) {祭司/史官};

% 分割线
\draw[white, thin, opacity=0.4]
  ({-\pw/2*0.85}, -4.2) -- ({\pw/2*0.85}, -4.2);

% 左侧标注
\draw[AccentGold!80, -{Stealth[length=5pt]}, thick]
  (-5.2, -5.0) -- (-3.0, -4.6);
\node[text width=3.2cm, align=right, font=\scriptsize\color{AccentGold!80},
      anchor=east] at (-5.3, -4.95)
  {\bfseries 生产平权\\（AI带来的变革）};

% 右侧标注：历史轴
\foreach \y/\era in {
  -0.75/公元前,
  -1.65/古代至近代,
  -2.60/20世纪,
  -3.65/互联网时代
}{
  \node[font=\scriptsize\color{WarmGray}, anchor=west] at (3.6, \y) {\era};
}

\end{tikzpicture}
\caption{信息生产权的历史金字塔：从古代祭司到今天的1.5\%创作者，
生产权始终是极少数人的特权——直到生成式AI出现}
\label{fig:pyramid}
\end{figure}

\vspace{6pt}

\subsection{真正的瓶颈：Know-how 的门槛}

各类创作工具不断提供模板，试图降低生产门槛——
但这只解决了"操作难度"的问题，并没有触及更深层的障碍：
\key{Know-how 的缺乏，以及表达能力本身的局限}。

举一个具体的例子：你想把昨夜的梦境还原出来，拍成一段影像。
这需要分镜脚本、摄影技术、剪辑能力、视觉特效……
任何一个环节的缺失，都会让你脑海中栩栩如生的画面，
永远停留在无法表达的沉默里。

\begin{blockquote}
「我想要一个什么样的东西」，这本身也是一种\textbf{人类表达}——
但在AI出现之前，绝大多数人的这种表达，
只能以「从市场上选择别人提供的供给」的方式勉强实现。
\end{blockquote}

生产表达被长期压抑，导致了大量个性化需求无从被满足。
人们不是不想创造，而是被客观能力和技术门槛挡在了门外。

% ============================================================
\section{通用人工智能：生产平权的来临}
% ============================================================

这正是生成式AI带来的\key{第一个结构性变化}：

它不只是一个更好用的工具，而是\key{将生产的权利还给了每一个人}。

\subsection{从选择到创造}

Vibe Coding 出现之前，一个普通人想开发一款自己的 App，
需要学习至少一门编程语言、理解系统架构、掌握调试方法——
这道门槛，拦住了绝大多数有想法的人。今天，
你只需要用语言描述清楚你想要什么，AI 可以帮你写出可运行的代码。

这不是说 AI 取代了工程师的全部价值，
而是说\emphasis{表达意图}本身，正在成为一种新的生产能力。

\begin{coreview}
信息消费侧的平权，已经由互联网和推荐算法基本完成。\\[4pt]
生成式AI所带来的，是\textbf{生产侧的平权}——\\
它让每一个普通人，都有能力将脑海中的想法，
转化为世界中真实存在的内容与产品。\\[4pt]
这是两次平权的闭环，也是信息革命的\textbf{终局意义}所在。
\end{coreview}

\subsection{「终局」的边界}

需要说明的是，这里所说的"终局"，
并非指技术进步从此终止，
而是指信息革命的\key{逻辑闭环}：
生产、分发、理解——三个环节都被AI渗透，
人类在信息链路上所面临的结构性门槛，
首次有了被系统性消解的可能。

当然，新的平权并不意味着新的绝对公平。
模型的算力、数据权利与平台规则，仍然可能形成新形式的集中。
但\key{个体把想法转化为产品的成本结构}，已经发生了根本性变化。

% ============================================================
\section{生产史的终极转向}
% ============================================================

\subsection{七个阶段：人类生产方式的演变}

要理解 AI 带来的这次变革有多深刻，
不妨将人类的生产历史拉通来看。

\vspace{14pt}

% ── 图3：人类生产历史演进 ────────────────────────────────────
\begin{figure}[h]
\centering
\begin{tikzpicture}[
  every node/.style={font=\scriptsize},
  stage/.style={
    rectangle, rounded corners=5pt,
    minimum width=1.55cm, minimum height=1.8cm,
    text width=1.45cm, align=center,
    draw=AccentBlue!30, fill=AccentBlue!7,
    line width=0.8pt
  },
  arrow/.style={-{Stealth[length=5pt]}, AccentBlue!50, thick},
  era label/.style={font=\tiny\color{WarmGray}, text width=1.45cm, align=center}
]

\def\sep{2.4}

% 七个阶段（各节点独立着色）
\node[stage, fill=AccentBlue!8,  draw=AccentBlue!40] (s0) at (0*\sep,0)
  {\textbf{采集狩猎}\\[3pt]\color{WarmGray}自然获取};
\node[stage, fill=AccentBlue!12, draw=AccentBlue!52] (s1) at (1*\sep,0)
  {\textbf{手工生产}\\[3pt]\color{WarmGray}个体手作};
\node[stage, fill=AccentBlue!16, draw=AccentBlue!64] (s2) at (2*\sep,0)
  {\textbf{蒸汽工厂}\\[3pt]\color{WarmGray}机器代劳};
\node[stage, fill=AccentBlue!20, draw=AccentBlue!74] (s3) at (3*\sep,0)
  {\textbf{电气分工}\\[3pt]\color{WarmGray}流水效率};
\node[stage, fill=AccentBlue!24, draw=AccentBlue!82] (s4) at (4*\sep,0)
  {\textbf{自动化}\\[3pt]\color{WarmGray}信息时代};
\node[stage, fill=AccentBlue!28, draw=AccentBlue!90] (s5) at (5*\sep,0)
  {\textbf{AI 智能}\\[3pt]\color{WarmGray}机器创造};
\node[stage, fill=AccentGold!12, draw=AccentGold!70] (s6) at (6*\sep,0)
  {\textbf{创造纪元}\\[3pt]\color{WarmGray}人类表达};

% 箭头
\foreach \i/\j in {0/1, 1/2, 2/3, 3/4, 4/5, 5/6}{
  \draw[arrow] (s\i.east) -- (s\j.west);
}

% 特殊标注最后一阶段
\draw[AccentGold, line width=1.5pt]
  (s6.north west) rectangle (s6.south east);

% 上方时代标注
\node[font=\tiny\bfseries\color{AccentBlue!60}, above=8pt of s2]
  {工业革命};
\node[font=\tiny\bfseries\color{AccentBlue!80}, above=8pt of s4]
  {信息革命};
\node[font=\tiny\bfseries\color{AccentGold!90}, above=8pt of s6]
  {\bfseries 智能革命};

% 底部轴注
\draw[TimelineGray, -{Stealth[length=4pt]}]
  (-0.9, -1.4) -- (14.6, -1.4)
  node[right, font=\tiny\color{WarmGray}]{历史进程};

\end{tikzpicture}
\caption{人类生产方式的七个阶段：从自然采集到创造纪元，每次跃迁都重新定义了
"劳动"的内涵与主体}
\label{fig:production}
\end{figure}

\vspace{6pt}

在这条演进曲线上，有一个耐人寻味的循环：
人类最早\emphasis{不从事生产}——他们直接从大自然获取食物和能源，
这是第一阶段。
而当 AI 可以承担几乎所有基础生活资料的生产之后，
人类将\emphasis{再次不从事生产}——这是第七阶段。

\subsection{两次"不生产"的本质差异}

同样是"不生产"，原因却截然不同。

\vspace{10pt}

% ── 图4：两次不生产的对比 ───────────────────────────────────
\begin{figure}[h]
\centering
\begin{tikzpicture}[
  every node/.style={font=\small},
  card/.style={
    rectangle, rounded corners=6pt,
    minimum width=5.8cm, minimum height=3.8cm,
    text width=5.4cm,
    draw, line width=0.8pt,
    align=left, inner sep=12pt
  }
]

% 左卡：原始社会
\node[card, fill=gray!6, draw=gray!35] (left) at (-3.5, 0) {
  \textbf{\color{WarmGray}原始社会的「不生产」}\\[6pt]
  \footnotesize
  \color{DarkText}
  \textbullet\ 原因：匮乏，无法生产\\[3pt]
  \textbullet\ 状态：只能采集，被动依赖\\[3pt]
  \textbullet\ 人类：生存资料不稳定\\[3pt]
  \textbullet\ 本质：\textbf{受困于物质不足}
};

% 右卡：AI时代
\node[card, fill=AccentGold!8, draw=AccentGold!50] (right) at (3.5, 0) {
  \textbf{\color{AccentGold}AI时代的「不生产」}\\[6pt]
  \footnotesize
  \color{DarkText}
  \textbullet\ 原因：充裕，无需生产\\[3pt]
  \textbullet\ 状态：主动创造，自我驱动\\[3pt]
  \textbullet\ 人类：专注于表达与意义\\[3pt]
  \textbullet\ 本质：\textbf{解放于物质充裕}
};

% 中间 VS
\node[circle, draw=AccentBlue!50, fill=white,
      minimum size=1.1cm, font=\bfseries\color{AccentBlue}] at (0,0) {VS};

% 底部说明
\node[font=\scriptsize, text=WarmGray, text width=11cm, align=center]
  at (0, -2.6)
  {两次「不生产」形式相似，但驱动逻辑完全相反——
  一个是被迫，一个是自由；一个是物质匮乏，一个是物质充裕之后的主动选择。};

\end{tikzpicture}
\caption{两次「不生产」的本质对比：相同的表象背后，是截然不同的文明逻辑}
\label{fig:compare}
\end{figure}

\vspace{4pt}

在 AI 智能生产时代完全成熟之后，
人类基础生活资料的生产将逐步由 AI 承担。
而人类将从事的，是另一种意义上的"工作"——
\key{创造、发明、探索，以及生活本身}。

这不是乌托邦式的懒人设想，而是一次劳动内涵的根本性升级：
人类的时间，将从消耗性的重复生产，
转向主动性的创造和对未知领域的探索。

% ============================================================
\section{结语：创造纪元的来临}
% ============================================================

综合以上两部分逻辑，可以得出一个清晰的判断：

\begin{coreview}
通用人工智能带给人类信息时代的，
是\textbf{分发平权之后的生产平权}。\\[6pt]
这种双重平权，将改变人类的生产组织方式、知识获取方式，
乃至对"工作"本身意义的理解。\\[6pt]
它不只是效率的提升，而是一次\textbf{文明范式的切换}。
\end{coreview}

\vspace{10pt}

回到麦克卢汉那句话——
「人类曾经以采集食物为生，而如今他们重新要以采集信息为生。」

这句话将要成为历史。

不是因为信息变得不再重要，而是因为\key{人类在信息链路中的位置将会改变}：
从被动的消费者和受限的生产者，
进化为能够\emphasis{自由表达、主动创造}的个体。

未来，人类以创造为生。
所有人都将投入精力去攻克未知的领域——
这是 AI 带来的，一个全新的人类纪元。

\vspace{30pt}
\begin{center}
  \color{AccentBlue!30}\rule{6cm}{0.5pt}
\end{center}

% ── 思考题 ──────────────────────────────────────────────────
\vspace{20pt}
\begin{thoughtbox}

\textbf{当生成式 AI 把「拍电影」的能力向所有人开放，}
\textbf{这个行业真正的稀缺会移向哪里？}

\vspace{10pt}
{\small\color{DarkText}
当视听内容的生产门槛持续下降，银幕上的供给不再稀缺，
真正的竞争或许会集中在另一件事上：
\key{谁能让一群人，愿意在同一时间、同一空间里，
把注意力交给你两个小时。}

\vspace{8pt}
如果注意力的动员才是核心稀缺，
影院是否会越来越多地向「组织者」而非「散客」收费？
粉丝应援场、品牌首映、团体包场——
这些今天看来边缘的玩法，
会不会成为未来院线收入结构中举足轻重的一极？

\vspace{12pt}
\textbf{更本质的问题是：}\\
在一个任何人都能创作的世界里，
\textbf{不是这片能不能做出来，而是别人为什么要来看？}\\
你会用什么机制——符号、社群、仪式感，还是别的什么——
来回答这个比拍摄更难的问题？
}

\end{thoughtbox}

\end{document}
